\chapter{The Architecture Design Space: Tensor Programs, Choices, and Ablations}\label{ch:archdesign}

\Cref{ch:gpt} built GPT-2 and \cref{ch:modern} built Qwen3. Between those two models sit a few
dozen design decisions: where the normalisation goes, which nonlinearity the MLP uses, how position
is encoded, how wide the MLP is relative to the residual stream, whether linear layers carry
biases, how weights are initialised, and which small tricks keep a large run from blowing up.
None of them changes what a transformer \emph{is}. Together they decide whether a 7B-parameter run
trains smoothly at a high learning rate or wastes a week of cluster time on loss spikes.

This chapter maps that design space. For each decision it gives the options, the mechanism that
makes one option better, the evidence the field has accumulated, and the cost in FLOPs, bytes and
kernels. It also shows how to \emph{run an ablation yourself}: a configurable model in
\code{llmsys/ablate.py}, a sweep driver, and a methodology for separating real effects from noise
on a small budget. Systems readers should watch for a recurring pattern: several ``modelling''
choices, such as dropping biases, RMSNorm, gated MLP widths that are multiples of 64, and head
dimensions of 64--128, are really systems choices made to suit the hardware.

\begin{figure}[t]
\centering
\resizebox{\linewidth}{!}{%
\begin{tikzpicture}[
  st/.style={draw=sysblue, fill=sysbluebg, rounded corners=2pt, minimum width=2.25cm,
    minimum height=1.05cm, align=center, font=\small, text width=2.1cm},
  num/.style={circle, fill=sysblue, text=white, font=\scriptsize\bfseries, inner sep=1.2pt},
  arr/.style={-{Stealth[length=2mm]}, thick, draw=black!60}]
\node[st] (a) at (0,0) {Tensor programs\\(einsum)};
\node[st] (b) at (2.75,0) {Layers \&\\initialisation};
\node[st] (c) at (5.5,0) {Normalisation\\placement};
\node[st] (d) at (8.25,0) {MLPs \&\\gating};
\node[st] (e) at (11,0) {Position\\encoding};
\node[st, fill=keygreenbg, draw=keygreen] (f) at (11,-1.75) {Hyper-\\parameters};
\node[st, fill=keygreenbg, draw=keygreen] (g) at (8.25,-1.75) {Stability\\tricks};
\node[st, fill=keygreenbg, draw=keygreen] (h) at (5.5,-1.75) {Block-level\\variants};
\node[st, fill=labpurplebg, draw=labpurple] (i) at (2.75,-1.75) {Ablations on\\a budget};
\node[st, fill=labpurplebg, draw=labpurple] (j) at (0,-1.75) {Lab 7: run\\the sweep};
\foreach \n/\k in {a/1,b/2,c/3,d/4,e/5,f/6,g/7,h/8,i/9,j/10}
  \node[num] at (\n.north west) {\k};
\draw[arr] (a)--(b); \draw[arr] (b)--(c); \draw[arr] (c)--(d); \draw[arr] (d)--(e);
\draw[arr] (e)--(f); \draw[arr] (f)--(g); \draw[arr] (g)--(h); \draw[arr] (h)--(i); \draw[arr] (i)--(j);
\end{tikzpicture}}
\caption{Roadmap of this chapter. The top row covers the components of a block, the bottom row
the global choices and the methodology for testing any of them.}
\label{fig:archdesign:roadmap}
\end{figure}

% =====================================================================
\section{Tensor programs with einsum and einops}\label{sec:archdesign:einsum}
\covers{CS336 L2; CS336 A1 §3.2}

Almost every operation in a transformer is either a \emph{contraction} (a sum of products over
some index, i.e.\ a matrix multiply in disguise), a \emph{rearrangement} of axes, or an elementwise
function. Index bookkeeping is where most implementation bugs hide: \code{x @ W.T} silently does
the wrong thing if \code{x} has its axes in the wrong order, and \code{transpose(-2, -1)} means
nothing to a reader without the shapes in mind. Einstein-summation notation fixes this by naming
every axis.

\paragraph{The rule.} An einsum string lists the axes of each input and of the output. Axes that
appear in an input but not in the output are summed over; axes that appear in the output are kept;
\code{...} stands for ``any leading axes, carried through unchanged''. A linear layer applied
to a batch of sequences is
\begin{equation}
y_{b,t,o} \;=\; \sum_{i} x_{b,t,i}\, W_{o,i}
\qquad\Longleftrightarrow\qquad
\texttt{einsum("...i,oi->...o", x, W)}.
\label{eq:archdesign:linear}
\end{equation}
The same notation expresses the attention score matrix, which contracts over the head dimension
and keeps batch, head, query and key positions:
\[
S_{b,h,q,k} = \sum_d Q_{b,h,q,d}\,K_{b,h,k,d}
\quad\Longleftrightarrow\quad
\texttt{einsum("bhqd,bhkd->bhqk", Q, K)} .
\]

\begin{figure}[t]
\centering
\begin{tikzpicture}[
  ax/.style={draw, rounded corners=1pt, minimum height=0.55cm, minimum width=0.7cm, font=\small\ttfamily},
  keep/.style={ax, fill=fdata!25, draw=keygreen},
  sum/.style={ax, fill=warnred!15, draw=warnred},
  lbl/.style={font=\scriptsize, text=black!70}]
\node[lbl] at (-1.3,0) {Q};
\node[keep] (q1) at (0,0) {b}; \node[keep] (q2) at (0.8,0) {h}; \node[keep] (q3) at (1.6,0) {q}; \node[sum] (q4) at (2.4,0) {d};
\node[lbl] at (-1.3,-1) {K};
\node[keep] (k1) at (0,-1) {b}; \node[keep] (k2) at (0.8,-1) {h}; \node[keep] (k3) at (1.6,-1) {k}; \node[sum] (k4) at (2.4,-1) {d};
\node[lbl] at (6.0,-0.5) {S};
\node[keep] (s1) at (7.0,-0.5) {b}; \node[keep] (s2) at (7.8,-0.5) {h}; \node[keep] (s3) at (8.6,-0.5) {q}; \node[keep] (s4) at (9.4,-0.5) {k};
\draw[-{Stealth}, keygreen, thick] (q3.north) to[out=40,in=140] (s3.north);
\draw[-{Stealth}, keygreen, thick] (k3.south) to[out=-30,in=-120] (s4.south);
\draw[warnred, thick, decorate, decoration={brace, amplitude=4pt, mirror}] (2.85,-1.3) -- (2.85,0.3);
\node[font=\scriptsize, text=warnred, align=left, anchor=west] at (3.0,-0.5) {d appears in both\\inputs, not in output:\\summed (contracted)};
\node[font=\scriptsize, text=keygreen, align=center] at (5.2,0.75) {b, h shared: batched};
\end{tikzpicture}
\caption{Reading \texttt{"bhqd,bhkd->bhqk"}. Axes shared by both inputs and kept in the output
(b, h) become batch dimensions of a batched GEMM; the axis that disappears (d) is the GEMM's
inner dimension; q and k become its rows and columns.}
\label{fig:archdesign:einsum}
\end{figure}

\paragraph{einops.} The \code{einops} library adds two companions with the same naming
discipline: \code{rearrange} for reshapes and transposes, and \code{reduce} for reductions. The
pattern \code{"b t (h d) -> b h t d"} splits the last axis into heads and moves them forward,
which replaces a \code{view} followed by a \code{transpose} whose correctness depends on the reader
remembering the layout.

\begin{lstlisting}[caption={Multi-head attention written with named axes (compare \cref{lst:archdesign:mha-torch}).},label={lst:archdesign:mha-einops}]
from einops import rearrange, einsum

def mha(x, Wq, Wk, Wv, Wo, n_heads):
    # x: (b, t, d_model); W*: (d_out, d_in)
    q = einsum(x, Wq, "b t i, o i -> b t o")
    k = einsum(x, Wk, "b t i, o i -> b t o")
    v = einsum(x, Wv, "b t i, o i -> b t o")
    q, k, v = (rearrange(z, "b t (h d) -> b h t d", h=n_heads) for z in (q, k, v))
    s = einsum(q, k, "b h q d, b h k d -> b h q k") / q.shape[-1] ** 0.5
    s = s.masked_fill(causal_mask(s.shape[-1]), float("-inf"))
    o = einsum(s.softmax(-1), v, "b h q k, b h k d -> b h q d")
    o = rearrange(o, "b h t d -> b t (h d)")
    return einsum(o, Wo, "b t i, o i -> b t o")
\end{lstlisting}

\begin{lstlisting}[caption={The same computation with positional axes, as in \texttt{llmsys/layers.py}.},label={lst:archdesign:mha-torch}]
def mha_torch(x, Wq, Wk, Wv, Wo, n_heads):
    B, T, D = x.shape
    split = lambda z: z.view(B, T, n_heads, D // n_heads).transpose(1, 2)
    q, k, v = split(x @ Wq.T), split(x @ Wk.T), split(x @ Wv.T)
    s = (q @ k.transpose(-2, -1)) / q.shape[-1] ** 0.5
    s = s.masked_fill(causal_mask(T), float("-inf"))
    o = (s.softmax(-1) @ v).transpose(1, 2).reshape(B, T, D)
    return o @ Wo.T
\end{lstlisting}

\paragraph{Memory order.} PyTorch tensors are row-major with explicit strides (\cref{ch:pytorch}).
\code{transpose} and \code{rearrange} patterns that only permute axes return \emph{views} with
permuted strides and copy nothing; the copy happens later, when an operation needs contiguous
memory (\code{reshape} on a non-contiguous view, or a kernel that requires it). That is why
\cref{lst:archdesign:mha-torch} uses \code{reshape} rather than \code{view} after the transpose.
Weights are stored as $(d_\text{out}, d_\text{in})$, the convention of \code{nn.Linear}, so a
forward pass computes $\mat{x}\mat{W}^\top$: each output feature is a contiguous row of
$\mat{W}$, which is also the layout that tensor-parallel column splits shard along
(\cref{ch:distributed}).

\begin{systemslens}[Einsum is a front end, not a kernel]
\code{torch.einsum} parses the string and lowers it to permutes, reshapes and a batched GEMM
(\code{bmm}), so it runs at the same speed as the hand-written version. Two things can still go
wrong. First, a permute that makes an operand non-contiguous can trigger a hidden copy; profile
(\cref{ch:perf}) rather than guess. Second, with three or more operands the \emph{contraction
order} matters: $(\mat{Q}\mat{K}^\top)\mat{V}$ costs $O(T^2 d)$ while $\mat{Q}(\mat{K}^\top\mat{V})$
costs $O(T d^2)$. PyTorch picks an order with \code{opt\_einsum} when it is installed, but for
attention the softmax sits between the two products and prevents reordering. Removing that
softmax to unlock the cheap order is exactly the idea behind linear attention (\cref{ch:attnalt}).
\end{systemslens}

% =====================================================================
\section{Layers from scratch: Linear, Embedding, and initialisation}\label{sec:archdesign:init}
\covers{CS336 A1 §3.3}

\Cref{lst:archdesign:linear} shows the two parameterised building blocks used by the ablation model
in \code{llmsys/ablate.py}, written without \code{nn.Linear} or \code{nn.Embedding} so that the
storage layout and the initialisation are explicit.

\begin{lstlisting}[caption={Linear and Embedding layers from scratch (\texttt{llmsys/ablate.py}).},label={lst:archdesign:linear}]
class Linear(nn.Module):
    def __init__(self, d_in, d_out, bias=False, std=None):
        super().__init__()
        self.weight = nn.Parameter(torch.empty(d_out, d_in))
        std = std if std is not None else math.sqrt(2.0 / (d_in + d_out))
        nn.init.trunc_normal_(self.weight, std=std, a=-3 * std, b=3 * std)
        self.bias = nn.Parameter(torch.zeros(d_out)) if bias else None

    def forward(self, x):                       # any leading axes
        y = torch.einsum("...i,oi->...o", x, self.weight)
        return y if self.bias is None else y + self.bias

class Embedding(nn.Module):
    def __init__(self, num, d, std=0.02):
        super().__init__()
        self.weight = nn.Parameter(torch.empty(num, d))
        nn.init.trunc_normal_(self.weight, std=std, a=-3 * std, b=3 * std)

    def forward(self, ids):                     # a gather, not a matmul
        return self.weight[ids]
\end{lstlisting}

\paragraph{Embeddings are gathers.} Mathematically an embedding is a one-hot vector times a
matrix, but nobody computes it that way: indexing reads $d$ numbers per token, while a one-hot matmul
would read the whole $V\times d$ table. The backward pass is a scatter-add into the rows that were
used, so the gradient is \emph{row-sparse}. Optimisers still update every row, because AdamW's
moments decay everywhere, which is one reason large-vocabulary embeddings are a noticeable share of
optimiser time.

\paragraph{Why initialisation matters.} Consider $y=\mat{W}x$ with $x$ having i.i.d.\ entries of
variance $\sigma_x^2$ and $\mat{W}$ having i.i.d.\ entries of variance $\sigma_w^2$. Each output is a
sum of $d_\text{in}$ independent products, so
\begin{equation}
\operatorname{Var}(y_o) = d_\text{in}\,\sigma_w^2\,\sigma_x^2 .
\label{eq:archdesign:var}
\end{equation}
Preserving the forward variance needs $\sigma_w^2 = 1/d_\text{in}$ (LeCun); preserving the
backward variance, where the roles swap, needs $1/d_\text{out}$. Glorot and Bengio (2010)
compromise with $\sigma_w^2 = 2/(d_\text{in}+d_\text{out})$, which is the default in
\cref{lst:archdesign:linear}. ReLU-family activations zero half the signal, which motivates He
initialisation, $2/d_\text{in}$. Drawing from a normal distribution \emph{truncated} at $\pm3\sigma$
removes the rare large weights that would otherwise create outlier features at initialisation; the
truncation shrinks the realised standard deviation by about 1.3\%, which is harmless.

\paragraph{Residual streams add variance.} A residual network adds each block's output to the
stream, $x_{\ell+1}=x_\ell+f_\ell(x_\ell)$. If each branch contributes variance comparable to the
stream, the stream's variance grows with depth. Pre-norm blocks normalise the branch \emph{input},
so each branch adds a roughly constant amount and the stream's RMS grows like $\sqrt{\ell}$. GPT-2
counteracts this by scaling the initial weights of each branch's output projection by
$1/\sqrt{2L}$ (two branches per block), which keeps the stream at unit scale at initialisation.
Without any normalisation the multiplicative coupling compounds and the stream grows
exponentially. \Cref{fig:archdesign:depth} shows all three regimes, measured on a random
48-block residual MLP stack.

\begin{figure}[t]
\centering
\begin{tikzpicture}
\begin{axis}[width=0.8\linewidth, height=5.2cm, xlabel={block index $\ell$},
  ylabel={RMS of residual stream}, ymode=log, xmin=0, xmax=48, grid=both,
  grid style={black!8}, legend pos=north west, legend style={font=\scriptsize},
  tick label style={font=\scriptsize}, label style={font=\small}]
\addplot[warnred, thick] coordinates {(0,1.004) (2,1.332) (4,1.767) (6,2.335) (8,3.066) (10,3.947) (12,5.018) (14,6.559) (16,8.395) (18,11.025) (20,14.105) (22,17.787) (24,22.676) (26,29.164) (28,39.087) (30,48.437) (32,63.397) (34,80.932) (36,104.676) (38,138.811) (40,174.512) (42,242.965) (44,313.951) (46,401.601) (48,500.979)};
\addlegendentry{no norm, fan-average init}
\addplot[sysblue, thick] coordinates {(0,1.001) (2,1.298) (4,1.536) (6,1.724) (8,1.886) (10,2.054) (12,2.194) (14,2.335) (16,2.442) (18,2.575) (20,2.696) (22,2.802) (24,2.949) (26,3.055) (28,3.149) (30,3.249) (32,3.330) (34,3.441) (36,3.512) (38,3.564) (40,3.626) (42,3.718) (44,3.792) (46,3.858) (48,3.950)};
\addlegendentry{pre-norm, fan-average init}
\addplot[keygreen, thick, dashed] coordinates {(0,1.001) (2,1.004) (4,1.007) (6,1.010) (8,1.013) (10,1.016) (12,1.019) (14,1.022) (16,1.027) (18,1.030) (20,1.033) (22,1.036) (24,1.038) (26,1.041) (28,1.045) (30,1.047) (32,1.050) (34,1.053) (36,1.056) (38,1.060) (40,1.064) (42,1.066) (44,1.070) (46,1.072) (48,1.077)};
\addlegendentry{pre-norm, output proj.\ $\times 1/\sqrt{2L}$}
\end{axis}
\end{tikzpicture}
\caption{Growth of the residual stream with depth at initialisation, for random inputs passed
through 48 residual blocks $x\leftarrow x+\mat{W}_2\,\mathrm{ReLU}(\mat{W}_1 N(x))$ with
$d=256$ (measured on CPU in this environment). With no norm the growth is exponential; pre-norm makes
it $\propto\sqrt{\ell}$; depth-scaled output projections keep it flat.}
\label{fig:archdesign:depth}
\end{figure}

\begin{napkinmath}[The loss at initialisation]
A correctly initialised language model should start out predicting a nearly uniform distribution, so
its cross-entropy at step 0 is close to $\ln V$: 10.82 for GPT-2's $V=50{,}257$, 11.93 for Qwen3's
$V=151{,}936$, and 5.55 for a byte-level model with $V=256$. If the logits have a small random spread
with standard deviation $\sigma_z$, the expected loss is about $\ln V+\sigma_z^2/2$. With the final
norm producing unit-RMS features of width $d$ and a head initialised with standard deviation $s$,
$\sigma_z\approx s\sqrt d$. For the $d=64$, $V=256$ ablation model with fan-average init,
$s=\sqrt{2/320}=0.079$, so $\sigma_z\approx0.63$ and the initial loss is about $5.55+0.2=5.75$, which
is what the model reports. An initial loss of 30 means the init is broken, not that the model is
``hard''. Checking this number is the first step of every debugging session
(\cref{sec:archdesign:ablate}).
\end{napkinmath}

\begin{deepdive}[Width-aware parameterisations]
Fan-in initialisation makes activations the right size at step 0, but it does not make the
\emph{updates} the right size as width grows: with Adam, a weight update has a roughly
width-independent per-entry magnitude, so its effect on activations grows with width. The maximal
update parameterisation ($\mu$P; Yang et al., 2021) rescales initialisation, learning rates and
output multipliers per layer type so that the optimal learning rate transfers from a narrow proxy to
a wide model. \Cref{ch:scalinglaws} uses this to tune hyperparameters cheaply.
\end{deepdive}

% =====================================================================
\section{Normalisation: what to normalise, and where}\label{sec:archdesign:norm}
\covers{CS336 L3; CS336 A1 §3.4.1}

\Cref{ch:gpt} implemented LayerNorm and \cref{sec:modern:rmsnorm} implemented RMSNorm. Here we
look at the two decisions that matter more than the formula: \emph{where} the norm sits in the
block, and whether there is one at all.

\paragraph{Placement.} \Cref{fig:archdesign:normplace} shows the four placements in use.
\begin{itemize}
\item \textbf{Post-norm} (the original transformer, BERT): $x\leftarrow N(x+f(x))$. Every
  residual path passes through a norm, so the identity path is not an identity. Xiong et
  al.~(2020) showed that this makes gradients near the output large relative to those of early layers
  at initialisation, which is why post-norm models need a learning-rate warmup and are fragile at
  depth.
\item \textbf{Pre-norm} (GPT-2 onwards, nearly every current LLM): $x\leftarrow x+f(N(x))$. The
  residual path is a clean identity from input to output, gradients flow to every layer unattenuated,
  and training tolerates larger learning rates. The price is the $\sqrt{L}$ stream growth of
  \cref{fig:archdesign:depth} and the need for a final norm before the LM head.
\item \textbf{Sandwich} (``double norm'' or peri-norm): $x\leftarrow x+N_b(f(N_a(x)))$. A second norm
  on the branch output bounds what each block can add to the stream. Gemma 2 normalises both
  before and after each sublayer, and OLMo 2 moved its norms onto the branch outputs; both are
  reported to improve stability. The norm still stays \emph{off} the residual path itself.
\item \textbf{Parallel}: $x\leftarrow x+\mathrm{Attn}(N(x))+\mathrm{MLP}(N(x))$ (GPT-J, PaLM). One
  norm serves both branches, and the first projections of attention and MLP can be fused into one
  GEMM; \cref{sec:archdesign:block} looks at the trade-off.
\end{itemize}

\begin{figure}[t]
\centering
\resizebox{\linewidth}{!}{%
\begin{tikzpicture}[
  op/.style={draw, rounded corners=2pt, minimum width=1.25cm, minimum height=0.5cm, font=\scriptsize},
  nrm/.style={op, fill=fmemory!25, draw=napkin},
  fn/.style={op, fill=fcompute!20, draw=sysblue},
  plus/.style={circle, draw, inner sep=0pt, minimum size=0.36cm, font=\scriptsize},
  arr/.style={-{Stealth[length=1.6mm]}, thick},
  ttl/.style={font=\small\bfseries}]
% post-norm
\begin{scope}[xshift=0cm]
\node[ttl] at (0.6,4.0) {post-norm};
\draw[arr] (0,-0.3) -- (0,0.8);
\node[fn] (f1) at (1.2,1.25) {$f$};
\draw[arr] (0,0.6) -| (f1.south);
\node[plus] (p1) at (0,1.9) {$+$};
\draw[arr] (0,0.8) -- (p1);
\draw[arr] (f1.north) |- (p1.east);
\node[nrm] (n1) at (0,2.6) {$N$};
\draw[arr] (p1) -- (n1);
\draw[arr] (n1) -- (0,3.55);
\node[font=\scriptsize, text=warnred, align=center] at (1.35,2.6) {norm on the\\residual path};
\end{scope}
% pre-norm
\begin{scope}[xshift=4.0cm]
\node[ttl] at (0.6,4.0) {pre-norm};
\draw[arr] (0,-0.3) -- (0,1.75);
\node[nrm] (n2) at (1.2,0.75) {$N$};
\node[fn] (f2) at (1.2,1.5) {$f$};
\draw[arr] (0,0.3) -| (n2.south);
\draw[arr] (n2) -- (f2);
\node[plus] (p2) at (0,2.3) {$+$};
\draw[arr] (f2.north) |- (p2.east);
\draw[arr] (p2) -- (0,3.55);
\node[font=\scriptsize, text=keygreen, align=center] at (-0.8,1.2) {clean\\identity};
\end{scope}
% sandwich
\begin{scope}[xshift=8.0cm]
\node[ttl] at (0.6,4.0) {sandwich};
\draw[arr] (0,-0.3) -- (0,2.4);
\node[nrm] (n3) at (1.2,0.5) {$N_a$};
\node[fn] (f3) at (1.2,1.2) {$f$};
\node[nrm] (n3b) at (1.2,1.9) {$N_b$};
\draw[arr] (0,0.1) -| (n3.south);
\draw[arr] (n3) -- (f3); \draw[arr] (f3) -- (n3b);
\node[plus] (p3) at (0,2.75) {$+$};
\draw[arr] (n3b.north) |- (p3.east);
\draw[arr] (p3) -- (0,3.55);
\end{scope}
% parallel
\begin{scope}[xshift=12.0cm]
\node[ttl] at (0.8,4.0) {parallel};
\draw[arr] (0,-0.3) -- (0,2.4);
\node[nrm] (n4) at (1.2,0.5) {$N$};
\node[fn] (fa) at (0.75,1.4) {attn};
\node[fn] (fm) at (2.05,1.4) {MLP};
\draw[arr] (0,0.1) -| (n4.south);
\draw[arr] (n4) -- (fa); \draw[arr] (n4) -- (fm);
\node[plus] (p4) at (0,2.75) {$+$};
\draw[arr] (fa.north) |- (p4.east);
\draw[arr] (fm.north) |- ([yshift=1pt]p4.east);
\draw[arr] (p4) -- (0,3.55);
\end{scope}
\end{tikzpicture}}
\caption{Where the norm goes, for one sublayer $f$ (attention or MLP); time flows upwards. Only
post-norm puts a normalisation \emph{on} the residual path.}
\label{fig:archdesign:normplace}
\end{figure}

\paragraph{No norm at all.} Removing normalisation (Lab~7, ``no norm'') leaves the model correct but
unstable: without the per-token rescaling, the stream grows as in \cref{fig:archdesign:depth}, the
attention logits grow with it, and the optimiser has to use a much smaller learning rate. Careful
initialisation schemes (Fixup, T-Fixup) and per-branch scalars can make norm-free transformers train,
but nobody trains large LLMs that way.

\paragraph{LayerNorm, RMSNorm, and dropping biases.} RMSNorm skips the mean subtraction and the bias.
In quality it matches LayerNorm in essentially every published comparison. The common justification is
``fewer FLOPs'', which is true but misleading, because norms are a vanishing fraction of a
transformer's FLOPs. The real gain is bytes moved: normalisation and other elementwise operators are
memory-bound, so their share of runtime is far larger than their share of FLOPs (Ivanov et al.,
2021). The same logic, plus some evidence of better stability, explains why most modern models also
drop the biases from linear layers.

\begin{napkinmath}[A norm is cheap in FLOPs and expensive in bytes]
Take one RMSNorm over a bf16 activation of $8192$ tokens $\times\,4096$ features on an H100. It
reads and writes $2\times 8192\times4096\times2\,\text{B}\approx134$\,MB and performs about 4 FLOPs
per element, roughly $1.3\times10^8$ FLOPs. At 3.35\,TB/s the memory traffic takes about 40\,$\mu$s;
at 989\,TFLOP/s the arithmetic would take 0.14\,$\mu$s. The arithmetic intensity is about 1 FLOP/byte,
nearly 300$\times$ below the H100's ridge point. For comparison, the $4096\times4096$ projection that
follows does $2\times8192\times4096^2\approx2.7\times10^{11}$ FLOPs in about 280\,$\mu$s at peak. So
an unfused norm adds roughly 15\% to the time of the matmul next to it, despite contributing 0.05\% of
the FLOPs. Fusing the norm into the neighbouring kernel (\cref{ch:perf}, \cref{app:triton}) removes
most of that cost.
\end{napkinmath}

% =====================================================================
\section{MLPs, activations, and gating}\label{sec:archdesign:mlp}
\covers{CS336 L3; CS336 A1 §3.4.2}

The position-wise MLP holds about two thirds of a dense transformer's parameters, so its form matters.
The classic form is $\mathrm{FFN}(x)=\mat{W}_2\,\sigma(\mat{W}_1x)$ with $d_\text{ff}=4d$.

\paragraph{Activations.} ReLU, $\max(0,z)$, was the original choice. GELU,
$z\,\Phi(z)$, is a smooth version that weights each input by the probability that it is positive under
a standard normal; GPT-2 uses its tanh approximation (\cref{ch:gpt}). SiLU (Swish),
$z\,\mathrm{sigmoid}(z)$, has the same shape with a slightly deeper negative dip. Squared ReLU,
$\max(0,z)^2$, appears in a few models (Primer; Nemotron-4 340B is reported to use it). \Cref{fig:archdesign:acts}
plots them.

\begin{figure}[t]
\centering
\begin{tikzpicture}
\begin{axis}[width=0.72\linewidth, height=5cm, xmin=-4, xmax=3, ymin=-0.6, ymax=3,
  xlabel={$z$}, ylabel={$\sigma(z)$}, grid=both, grid style={black!8}, samples=120,
  legend pos=north west, legend style={font=\scriptsize}, tick label style={font=\scriptsize}]
\addplot[black!60, thick, domain=-4:3] {max(0,x)}; \addlegendentry{ReLU}
\addplot[sysblue, thick, domain=-4:3] {0.5*x*(1+tanh(0.7978845608*(x+0.044715*x^3)))}; \addlegendentry{GELU (tanh approx.)}
\addplot[keygreen, thick, dashed, domain=-4:3] {x/(1+exp(-x))}; \addlegendentry{SiLU / Swish}
\addplot[warnred, thick, dotted, domain=-4:1.73] {max(0,x)^2}; \addlegendentry{squared ReLU}
\end{axis}
\end{tikzpicture}
\caption{The common MLP nonlinearities. GELU and SiLU are smooth and non-monotonic near zero; the
gated variants of \cref{eq:archdesign:glu} multiply one of these curves by a second linear projection.}
\label{fig:archdesign:acts}
\end{figure}

\paragraph{Gated linear units.} A GLU MLP replaces $\sigma(\mat{W}_1x)$ with the elementwise product
of an activated projection and a linear one:
\begin{equation}
\mathrm{FFN}_\text{GLU}(x) = \mat{W}_2\big(\sigma(\mat{W}_g x)\odot \mat{W}_1 x\big),
\label{eq:archdesign:glu}
\end{equation}
which gives ReGLU ($\sigma=$ ReLU), GeGLU (GELU) and SwiGLU (SiLU). The extra matrix costs parameters,
so GLU MLPs use $d_\text{ff}\approx\frac{8}{3}d$, which makes $3\cdot\frac83 d^2 = 8d^2$ equal to the
$2\cdot4d^2$ of a plain MLP. Shazeer (2020) found that GLU variants consistently improved perplexity at
matched parameters, and Narang et al.~(2021) found GLUs among the few modifications whose gains
transferred across settings. Almost every LLM released since 2023 uses SwiGLU or GeGLU.

\begin{lstlisting}[caption={A parameter-matched MLP family (\texttt{llmsys/ablate.py}).},label={lst:archdesign:mlp}]
GATED = {"swiglu", "geglu", "reglu"}

def ffn_hidden_dim(d_model, kind, multiple_of=64):
    raw = 8 * d_model / 3 if kind in GATED else 4 * d_model
    return int(multiple_of * math.ceil(raw / multiple_of))   # tensor-core friendly

class MLP(nn.Module):
    def __init__(self, cfg, out_std):
        super().__init__()
        d, h = cfg.d_model, cfg.ff_dim()
        self.gated = cfg.ffn in GATED
        self.act = {"swiglu": F.silu, "geglu": F.gelu, "reglu": F.relu}.get(cfg.ffn) \
            or ACTS[cfg.ffn]
        self.w_in = Linear(d, h)
        self.w_gate = Linear(d, h) if self.gated else None
        self.w_out = Linear(h, d, std=out_std)        # depth-scaled init

    def forward(self, x):
        if self.gated:
            return self.w_out(self.act(self.w_gate(x)) * self.w_in(x))
        return self.w_out(self.act(self.w_in(x)))
\end{lstlisting}

\begin{systemslens}[What gating costs on the GPU]
A SwiGLU MLP at matched parameters does the same matmul FLOPs as a 4$d$ MLP, but in \emph{three} GEMMs
of width $\frac83 d$ instead of two of width $4d$. Production implementations concatenate
$\mat{W}_g$ and $\mat{W}_1$ into one $(2d_\text{ff})\times d$ weight (the ``\code{gate\_up\_proj}'' in
many inference engines), so the hardware sees one wide GEMM, then a fused
\code{silu(a) * b} elementwise kernel. For the backward pass, a SwiGLU MLP saves three hidden-width
tensors per token (gate pre-activation, up projection, and their product), while a GELU MLP saves two
(pre-activation and activation). Measured in $d_\text{ff}$ that sounds worse, but at matched parameters
$3\cdot\frac83d=8d=2\cdot4d$, so the activation memory is the same; an unfused implementation, however,
also materialises $\mathrm{SiLU}(\vect{a})$ separately, which is exactly what fusion removes
(\cref{ch:perf} shows how to measure this). Finally, $\frac83 d$ is rarely a nice number. Rounding it up to a multiple of 64 or 128 keeps every GEMM tile full; Qwen3-0.6B's
$d_\text{ff}=3072=3d$ is one such rounding choice.
\end{systemslens}

% =====================================================================
\section{Positional information}\label{sec:archdesign:pos}
\covers{CS336 L3; CS336 A1 §3.4.3, §7.3}

Self-attention is permutation-equivariant: without extra information, swapping two input tokens just swaps
the corresponding outputs. Something has to tell the model where each token sits. Five families are in use.

\paragraph{No positional encoding (NoPE).} A \emph{causal} decoder is not actually permutation-invariant,
because the mask itself breaks the symmetry. A token at position $i$ attends over exactly $i+1$ keys, so
quantities such as the attention entropy, or the average of a constant value vector, already depend on
$i$. Theory (Kazemnejad et al., 2023) and experiments show that causal transformers trained without any
explicit encoding can learn both absolute and relative positional information. NoPE is also used
deliberately in modern hybrids, for example in the global-attention layers of some models that use RoPE
only in their local layers.

\paragraph{Absolute embeddings.} GPT-2 adds a learned vector $\vect{u}_i$ per position
(\cref{ch:text}); the original transformer added fixed sinusoids,
\[
\mathrm{PE}_{i,2j}=\sin(i/10000^{2j/d}),\qquad \mathrm{PE}_{i,2j+1}=\cos(i/10000^{2j/d}).
\]
Sinusoids have a relative structure: $\langle\mathrm{PE}_i,\mathrm{PE}_{i+\Delta}\rangle$ depends only on
$\Delta$. But once they are \emph{added} to content embeddings, the query--key product
$\langle\vect{x}+\mathrm{PE}_i,\vect{y}+\mathrm{PE}_j\rangle$ contains cross terms such as
$\langle\vect{x},\mathrm{PE}_j\rangle$ that depend on absolute positions. Learned absolute embeddings
cannot extrapolate past the trained length at all, because those rows were never trained.

\paragraph{Relative biases and ALiBi.} T5 adds a learned scalar bias to each attention logit, indexed by a
bucketed distance $i-j$. ALiBi (Press et al., 2022) uses a fixed bias that grows linearly with distance,
\[
s_{ij} \leftarrow s_{ij} - m_h\,(i-j),\qquad m_h = 2^{-8h/H},
\]
with a geometric set of slopes so that some heads are local and others nearly global. It adds no parameters
and degrades gracefully beyond the training length.

\paragraph{RoPE.} Rotary embeddings (\cref{sec:modern:rope}) satisfy the property we actually want: a
scoring function of the form $\langle f(\vect{x},i), f(\vect{y},j)\rangle = g(\vect{x},\vect{y},i-j)$.
Rotating queries and keys by position-dependent angles achieves this exactly, because rotations
preserve inner products, and there are no additive cross terms. RoPE is the default for essentially
all current open LLMs; partial RoPE (rotating only some dimensions) and per-layer variations are
common refinements.

\begin{figure}[t]
\centering
\begin{tikzpicture}
\begin{axis}[width=0.52\linewidth, height=4.6cm, xmin=0, xmax=64, ymin=-1.15, ymax=1.15,
  xlabel={position $i$}, ylabel={$\mathrm{PE}_{i,2j}$}, title={\small sinusoidal channels ($d=64$)},
  samples=160, tick label style={font=\scriptsize}, legend style={font=\tiny, at={(0.98,0.02)}, anchor=south east}]
\addplot[sysblue, thick, domain=0:64] {sin(deg(x))}; \addlegendentry{$j=0$}
\addplot[keygreen, thick, domain=0:64] {sin(deg(x/10000^(8/64)))}; \addlegendentry{$j=4$}
\addplot[warnred, thick, domain=0:64] {sin(deg(x/10000^(24/64)))}; \addlegendentry{$j=12$}
\end{axis}
\begin{scope}[xshift=8.1cm, yshift=0.3cm, x=0.42cm, y=0.42cm]
\node[font=\small] at (4,8.6) {ALiBi bias, one head ($m=\tfrac12$)};
\foreach \i in {0,...,7} {
  \foreach \j in {0,...,7} {
    \ifnum\j>\i
      \fill[black!4] (\j,7-\i) rectangle ++(1,1);
    \else
      \pgfmathsetmacro{\shade}{100-12*(\i-\j)}
      \fill[sysblue!\shade!white] (\j,7-\i) rectangle ++(1,1);
    \fi
  }
}
\draw[black!40] (0,0) grid (8,8);
\node[font=\scriptsize, rotate=90] at (-0.6,4) {query $i$};
\node[font=\scriptsize] at (4,-0.6) {key $j$};
\node[font=\tiny, text=white] at (0.5,0.5) {$-3.5$};
\node[font=\tiny] at (7.5,0.5) {$0$};
\end{scope}
\end{tikzpicture}
\caption{Left: sinusoidal encodings use a geometric range of frequencies, so low-index channels change
quickly with position and high-index channels slowly. Right: ALiBi subtracts $m(i-j)$ from each logit;
darker cells are closer keys, which receive less penalty. Cells above the diagonal are masked.}
\label{fig:archdesign:pos}
\end{figure}

\begin{systemslens}[Positional schemes and the inference stack]
The choice of positional scheme reaches into the inference engine. Absolute embeddings are added once at the
input and never touch attention. RoPE is applied to $\vect{q}$ and $\vect{k}$ \emph{before} keys are
cached, so the cache stores already-rotated keys: prefix caching (\cref{ch:serving}) stays valid only if a
shared prefix sits at the same positions, and positions shifted by left padding are harmless because only
relative angles matter (\cref{ch:inference}). ALiBi and T5-style biases must be added \emph{inside} the
attention kernel, which FlashAttention supports, but each such bias is one more variant that kernels must
specialise. NoPE costs nothing anywhere, which is part of its appeal for long-context global layers.
\end{systemslens}

% =====================================================================
\section{Hyperparameters that (mostly) do not need tuning}\label{sec:archdesign:hparams}
\covers{CS336 L3}

A transformer's shape is fixed by a handful of numbers: $d$, $L$, $H$, $\dhead$, $d_\text{ff}$ and $V$.
Across many published models these have settled into narrow ranges (\cref{tab:archdesign:shapes}).

\begin{table}[t]
\centering\small
\begin{tabular}{@{}lrrrrrrr@{}}
\toprule
Model & $d$ & $L$ & $H$ (KV) & $\dhead$ & $d_\text{ff}$ & $d_\text{ff}/d$ & $d/L$\\
\midrule
GPT-2 small & 768 & 12 & 12 & 64 & 3072 & 4.0 & 64\\
GPT-2 XL & 1600 & 48 & 25 & 64 & 6400 & 4.0 & 33\\
GPT-3 175B & 12288 & 96 & 96 & 128 & 49152 & 4.0 & 128\\
Llama 2 70B & 8192 & 80 & 64 (8) & 128 & 28672 & 3.5 & 102\\
Qwen3-0.6B & 1024 & 28 & 16 (8) & 128 & 3072 & 3.0 & 37\\
Qwen3-8B & 4096 & 36 & 32 (8) & 128 & 12288 & 3.0 & 114\\
Qwen3-32B & 5120 & 64 & 64 (8) & 128 & 25600 & 5.0 & 80\\
\bottomrule
\end{tabular}
\caption{Shapes of some well-documented dense models. Qwen3 and Llama use gated MLPs, so the
parameter-matched ratio is $\frac83\approx2.67$; GPT-2 and GPT-3 use plain 4$d$ MLPs. In Qwen3-0.6B and
Qwen3-32B, $H\cdot\dhead\neq d$: the attention width is decoupled from the residual width.}
\label{tab:archdesign:shapes}
\end{table}

\paragraph{MLP ratio.} $d_\text{ff}=4d$ for plain MLPs and about $\frac83 d$ to $3.5d$ for GLUs. Kaplan et
al.~(2020) found loss to be flat across a wide range of ratios, and T5's 11B model used $d_\text{ff}=64d$
without disaster. The default is a safe point in a broad basin, not a sharp optimum.

\paragraph{Heads and head dimension.} The convention $H\cdot\dhead=d$ keeps attention's projections square,
but nothing requires it, as two of the Qwen3 configurations show. What matters for systems is
$\dhead$: the FlashAttention tile budget in \cref{sec:attention:variants} makes 64--128 the efficient range
on current GPUs, and kernels are specialised per head dimension. The number of KV heads is an
inference decision (\cref{sec:modern:gqa}).

\paragraph{Aspect ratio.} Wide-and-shallow and narrow-and-deep models of the same parameter count reach
similar loss over a broad range of $d/L$, roughly 50 to 200 (Kaplan et al., 2020). Systems
considerations usually decide the value: depth adds sequential latency per decoded token and more pipeline
stages, while width makes each GEMM larger and easier to split with tensor parallelism.

\paragraph{Vocabulary size.} Monolingual English models use 32k--50k tokens; multilingual and production
models use 100k--260k (\cref{ch:text}). Larger vocabularies shorten sequences but enlarge the embedding
table and the output softmax, which dominate the parameter count of small models: in Qwen3-0.6B the tied
$151{,}936\times1024$ embedding is 156M of its roughly 600M parameters.

\paragraph{Regularisation.} Modern pretraining runs are single-epoch, so overfitting is not the concern.
Most recent models use no dropout and weight decay of about 0.1. Andriushchenko et al.~(2023) argue that
weight decay in LLM training acts mainly on optimisation dynamics, through its interaction with the
learning-rate schedule, rather than as a regulariser. \Cref{ch:optim} implements decoupled weight decay.

\begin{napkinmath}[Pad your shapes]
Tensor-core GEMMs work in tiles; on an H100 bf16 kernels are happiest when the matrix dimensions are
multiples of 64 or 128. GPT-2's $V=50{,}257$ is prime-ish, so the LM-head GEMM ($T\times d\times V$) has a
ragged last tile. Padding the vocabulary to $50{,}304=64\times786$ adds 47 unused rows (0.09\% more
parameters) and makes every tile full. In small-model training, where the LM head is a large share of the
FLOPs, practitioners report noticeable speedups from this one change. The same reasoning applies to
$d_\text{ff}$ (round $\frac83d$ up to a multiple of 64) and to batch $\times$ sequence products.
\end{napkinmath}

% =====================================================================
\section{Stability tricks}\label{sec:archdesign:stability}
\covers{CS336 L3}

Large runs fail in a characteristic way: the loss curve develops spikes, some of which recover and some
of which do not. Most of the fixes target the two softmaxes in the model, because exponentials turn
moderate growth in logits into overflow and saturated gradients.

\paragraph{Output softmax: the z-loss.} The cross-entropy depends only on logit \emph{differences}, so
nothing in the loss stops all logits from drifting upward together. The log-partition function
$\log Z=\operatorname{logsumexp}(\vect{z})$ can wander far from zero, which in bf16 costs precision and
can eventually overflow. PaLM added an auxiliary term
\begin{equation}
\mathcal{L} = \mathcal{L}_\text{CE} + \alpha\,(\log Z)^2,\qquad \alpha\approx10^{-4},
\label{eq:archdesign:zloss}
\end{equation}
whose gradient, $2\alpha\log Z\cdot\softmax(\vect{z})$, gently pulls $\log Z$ back toward 0 without changing
the predicted distribution. Several later open models (Baichuan 2, DCLM, OLMo 2) also use it.

\paragraph{Attention softmax: QK-norm.} As training proceeds, attention logits $\vect{q}^\top\vect{k}/\sqrt{\dhead}$
can grow until the softmax becomes one-hot. Gradients through it then vanish, which is a known cause of
divergence (Dehghani et al., 2023). Normalising $\vect{q}$ and $\vect{k}$ per head before the dot product
bounds the logits by the learned norm scales. Qwen3 adopted QK-norm (\cref{sec:modern:gqa}), as did OLMo 2
and Gemma 3, among others.

\paragraph{Logit soft-capping.} Gemma 2 instead bounds logits smoothly with $c\tanh(z/c)$
(\cref{fig:archdesign:softcap}): nearly the identity for $|z|\ll c$, saturating at $\pm c$. It caps
attention logits at 50 and final logits at 30. The tanh does not fit the standard fused attention
kernels without modification, which is one reason QK-norm has been the more popular route.

\begin{figure}[t]
\centering
\begin{tikzpicture}
\begin{axis}[width=0.6\linewidth, height=4.6cm, xmin=-120, xmax=120, ymin=-60, ymax=60,
  xlabel={raw logit $z$}, ylabel={capped logit}, grid=both, grid style={black!8}, samples=150,
  legend pos=north west, legend style={font=\scriptsize}, tick label style={font=\scriptsize}]
\addplot[black!40, dashed, domain=-60:60] {x}; \addlegendentry{identity}
\addplot[sysblue, thick, domain=-120:120] {30*tanh(x/30)}; \addlegendentry{$30\tanh(z/30)$}
\addplot[keygreen, thick, domain=-120:120] {50*tanh(x/50)}; \addlegendentry{$50\tanh(z/50)$}
\end{axis}
\end{tikzpicture}
\caption{Logit soft-capping. Ordinary logits pass through almost unchanged; runaway logits saturate
smoothly instead of overflowing, and gradients never become exactly zero.}
\label{fig:archdesign:softcap}
\end{figure}

\begin{lstlisting}[caption={Cross-entropy with z-loss, computed in fp32 (\texttt{llmsys/ablate.py}).},label={lst:archdesign:zloss}]
def lm_loss(logits, targets, z_loss=0.0):
    lf = logits.float().flatten(0, -2)                 # (N, V) in fp32
    lse = torch.logsumexp(lf, dim=-1)                  # log Z per position
    ce = (lse - lf.gather(-1, targets.reshape(-1, 1)).squeeze(-1)).mean()
    total = ce + z_loss * lse.pow(2).mean() if z_loss > 0 else ce
    return total, ce.detach(), lse.mean().detach()     # log the size of log Z too
\end{lstlisting}

\paragraph{Everything else.} The remaining items on the stability checklist are less glamorous:
\begin{itemize}
\item compute softmaxes, norms and the loss in fp32 even when matmuls run in bf16 (\cref{ch:pretrain});
\item use depth-scaled initialisation (\cref{fig:archdesign:depth}) and no biases;
\item use warmup, gradient clipping and $\beta_2=0.95$ rather than 0.999 for AdamW (\cref{ch:optim});
\item in MoE models, run the router in fp32 and add a router z-loss (\cref{ch:attnalt});
\item log the gradient norm, the mean $\log Z$, the maximum attention logit and the per-layer activation
  RMS. Each of these usually starts drifting several hundred steps before a spike, which makes them
  early warnings.
\end{itemize}

\begin{systemslens}[A spike costs more than the steps it wastes]
At scale a divergence is a systems event. The usual recovery is to roll back to a checkpoint from before
the spike, skip or reshuffle the data batches that triggered it, and sometimes lower the learning rate.
That requires frequent asynchronous checkpoints (\cref{ch:pretrain}), deterministic data loading that can
skip specific batches, and monitoring that alerts within minutes. Several published training logs (OPT,
BLOOM) describe dozens of manual restarts. Stability tricks are cheap insurance against an expensive
operational failure mode.
\end{systemslens}

% =====================================================================
\section{Block-level variants}\label{sec:archdesign:block}
\covers{CS336 L3}

\paragraph{Serial versus parallel blocks.} A parallel block computes attention and MLP from the same
normalised input and adds both to the stream. Because both branches start with a matmul of the same
input, their weights can be concatenated into one $(3d + 2d_\text{ff})\times d$ GEMM, which improves
utilisation at small batch sizes and lets tensor parallelism overlap one all-reduce for both branches.
GPT-J, GPT-NeoX and PaLM used parallel blocks, and PaLM reported about 15\% faster training at large
scale with a small quality cost at smaller sizes. Most recent models have returned to serial blocks.

\paragraph{Biases.} Dropping all bias vectors removes parameters that the optimiser must still track
in fp32 (master weights plus two Adam moments), removes elementwise kernels or epilogues, and has been
reported to help stability. The cost in quality has not been measurable at LLM scale, with one common
exception: Qwen2 kept a bias on the QKV projections.

\paragraph{Tied embeddings.} Sharing the input embedding and the output head saves $Vd$ parameters, which
is a large fraction for small models (GPT-2, Qwen3 up to 4B) and negligible for large ones (Qwen3-8B and
larger untie them). Tying couples two roles with different gradient statistics, so large models prefer
the flexibility.

% =====================================================================
\section{What recent open models chose}\label{sec:archdesign:survey}
\covers{CS336 L3}

\Cref{tab:archdesign:survey} summarises the architectural choices of several widely studied models, as
reported in their papers and released configurations. The convergence is striking: pre-norm (often with
extra branch norms), RMSNorm, gated MLPs, RoPE and no biases make up the ``Llama-like'' default, and new
models mostly vary the attention pattern (GQA, local/global interleaving, hybrids;
\cref{ch:attnalt}) and the stability tricks.

\begin{table}[t]
\centering\small
\begin{tabularx}{\linewidth}{@{}l>{\raggedright\arraybackslash}X>{\raggedright\arraybackslash}X>{\raggedright\arraybackslash}X>{\raggedright\arraybackslash}X@{}}
\toprule
Model & Norm & MLP & Position & Notable extras\\
\midrule
GPT-2 & pre-LN, LayerNorm & GELU, 4$d$ & learned absolute & biases; tied embeddings\\
GPT-3 & pre-LN, LayerNorm & GELU, 4$d$ & learned absolute & alternating dense and locally banded sparse attention\\
PaLM & pre-norm, LayerNorm without bias & SwiGLU & RoPE & parallel blocks; MQA; z-loss; no biases\\
Llama 2 & pre-RMSNorm & SwiGLU & RoPE & GQA in the 70B model; no biases\\
Gemma 2 & RMSNorm before and after each sublayer & GeGLU & RoPE & logit soft-capping; alternating local/global attention\\
OLMo 2 & RMSNorm on branch outputs & SwiGLU & RoPE & QK-norm; z-loss\\
Qwen3 (dense) & pre-RMSNorm & SwiGLU & RoPE, $\theta=10^6$ & QK-norm; GQA; no biases\\
\bottomrule
\end{tabularx}
\caption{Architectural choices of some well-documented models, as reported by their authors. Details
vary across sizes within a family; always check the released config of the specific checkpoint.}
\label{tab:archdesign:survey}
\end{table}

% =====================================================================
\section{Running ablations on a budget}\label{sec:archdesign:ablate}
\covers{CS336 A1 §7.2.3, §7.3}

Everything above is a claim made by somebody else. The only way to build real intuition, and to settle
questions for your own setting, is to run controlled experiments. Small-scale ablations are cheap,
but they are full of traps. This section covers the setup, the methodology and the debugging
habits, and then reports a sweep you can reproduce in a few CPU-minutes.

\subsection{The setup: one config, one model, one loop}

The ablation model in \code{llmsys/ablate.py} exposes every choice in this chapter as a field of a
dataclass, so a variant is a one-line override and the full configuration is logged with the result:

\begin{lstlisting}[caption={Configuring and running ablations.},label={lst:archdesign:config}]
from dataclasses import replace
from llmsys.ablate import ArchConfig, train_lm, make_corpus

base = ArchConfig(d_model=64, n_layers=4, n_heads=4, context_length=64,
                  norm="rms", norm_position="pre", pos="rope", ffn="swiglu")
ids = make_corpus(300_000)            # or: bytes of a real text file
variants = {
    "baseline": {},
    "post-norm": {"norm_position": "post"},
    "no norm": {"norm": "none"},
    "NoPE": {"pos": "nope"},
    "SiLU MLP (4d)": {"ffn": "silu"},        # parameter-matched to SwiGLU
}
results = [train_lm(replace(base, **ov), ids, steps=250, lr=3e-3, seed=s, name=n)
           for n, ov in variants.items() for s in (0, 1)]
\end{lstlisting}

\code{train\_lm} is a deliberately complete loop: AdamW with $\beta_2=0.95$ and decoupled weight decay on
matrices only, linear warmup and cosine decay, gradient clipping, periodic validation on a held-out
slice, gradient-norm and $\log Z$ logging, and divergence detection that records a NaN as a
\emph{result} rather than crashing the sweep. \code{labs/lab07\_ablations/run.py} wraps it with a
command-line interface, saves JSON after every run so an interrupted sweep loses at most one run, and
plots the curves. On a GPU, point \code{--text} at a real corpus (for example a TinyStories dump) and
scale the shape up.

\subsection{Methodology}

\begin{enumerate}
\item \textbf{Change one thing.} Each variant differs from the baseline in exactly one field, and all
  variants share the data order (same seed for the data sampler) so that differences are not caused by
  which batches happened to be drawn.
\item \textbf{Match what should be matched.} Compare at equal parameters (hence $4d$ for non-gated MLPs)
  and equal tokens. When a change alters FLOPs per token, also compare at equal FLOPs or equal
  wall-clock time.
\item \textbf{Tune the learning rate per variant}, at least coarsely. Many ``X is better than Y'' results
  are really ``X tolerates the learning rate that was tuned for X''. Post-norm and no-norm models in
  particular need lower rates, so a single-LR comparison mostly measures LR sensitivity. Report
  the LR sweep too.
\item \textbf{Know your noise floor.} Run the baseline with two or three seeds. Differences smaller than
  the seed-to-seed spread are not results.
\item \textbf{Look at curves, not just endpoints.} A variant that starts slowly may cross over later,
  and early curves can mislead in both directions.
\item \textbf{Be humble about transfer.} Narang et al.~(2021) re-tested dozens of published transformer
  modifications and found that most did not transfer across codebases and scales. A small-scale win is a
  hypothesis to test at the next scale up, not a conclusion.
\end{enumerate}

\begin{pitfall}[Debugging before ablating]
Ablations on a broken model measure nothing. Before any sweep, check, in order:
(1) the loss at step 0 is close to $\ln V$ (\cref{sec:archdesign:init});
(2) the model can \emph{overfit a single batch} to near-zero loss within a few hundred steps, since if it
cannot, there is a bug in the model, the loss or the optimiser;
(3) tensor shapes match your expectations at every boundary (assert them, or name them with einops);
(4) the norms of activations, weights and gradients are finite and roughly constant across layers.
Step through the forward pass in a debugger rather than adding prints. Most ``my architecture
change hurt'' results at small scale turn out to be a mask, a transpose or an init bug.
\end{pitfall}

\subsection{A worked sweep}
%ABLATION-RESULTS%

% =====================================================================
\section{Summary}\label{sec:archdesign:summary}
\begin{itemize}
\item Einsum and einops name every axis; they lower to the same batched GEMMs as hand-written code, but
  remove whole classes of transpose bugs. Weights are stored $(d_\text{out},d_\text{in})$.
\item Initialisation is variance bookkeeping: fan-average truncated normals for matrices, small embeddings,
  and $1/\sqrt{2L}$ scaling of residual output projections. The loss at step 0 should be about $\ln V$.
\item Pre-norm keeps the residual path an identity and is why deep transformers train at high learning
  rates; extra norms on branch outputs are a common refinement. RMSNorm and bias-free layers win mainly
  by moving fewer bytes.
\item Gated MLPs (SwiGLU, GeGLU) at $d_\text{ff}\approx\frac83 d$ are the consensus; implement them as one
  fused gate-up GEMM.
\item RoPE dominates positional encoding. NoPE works in causal models, and ALiBi and relative biases live
  inside the attention kernel.
\item Shapes sit in broad basins ($d_\text{ff}/d\approx 3$--4, $d/L\approx 50$--200, $\dhead=64$--128);
  systems constraints usually decide the exact values.
\item Stability tricks target softmaxes: z-loss for the output and QK-norm or soft-capping for
  attention, plus fp32 reductions and monitoring.
\item Ablations need one change at a time, matched parameters and tokens, per-variant learning rates, a
  noise floor, and a debugged baseline.
\end{itemize}

% =====================================================================
\section{Exercises}\label{sec:archdesign:exercises}

\begin{exercise}[einsum fluency]
Write einsum strings for: (a) the GQA score computation where queries have shape $(b,g,r,q,d)$ with
$g$ KV groups and $r$ queries per group, and keys $(b,g,k,d)$; (b) a batched outer product
$(b,i),(b,j)\to(b,i,j)$; (c) the trace of each matrix in a batch $(b,n,n)$; (d) the sum over heads of
per-head output projections, given $\mat{O}\in\R^{b\times h\times t\times d_h}$ and
$\mat{W}\in\R^{h\times d_h\times D}$.
\end{exercise}

\begin{exercise}[Contraction order]
For $\mat{A}\in\R^{n\times d}$, $\mat{B}\in\R^{d\times n}$, $\mat{C}\in\R^{n\times d}$, count the FLOPs of
$(\mat{A}\mat{B})\mat{C}$ and $\mat{A}(\mat{B}\mat{C})$. For $d=128$, at what $n$ does the second order
become cheaper, and by how much at $n=32{,}768$? Which nonlinearity in attention prevents using the
cheaper order?
\end{exercise}

\begin{exercise}[Variance propagation]
Derive the variance of the output of a 2-layer ReLU MLP $y=\mat{W}_2\,\mathrm{ReLU}(\mat{W}_1x)$ under
fan-average initialisation, for $x\sim\mathcal{N}(0,\mat{I}_d)$ and $d_\text{ff}=4d$. What scale on
$\mat{W}_2$ keeps $\operatorname{Var}(y)=\operatorname{Var}(x)$?
\end{exercise}

\begin{exercise}[Loss at init]
Predict the initial loss of (a) Qwen3-0.6B with a head initialised with standard deviation 0.02 and
unit-RMS final features; (b) the same model with the head initialised at standard deviation 1. Verify
(a) with the core \code{Qwen3Model} on random tokens (no download needed).
\end{exercise}

\begin{exercise}[Parameter matching]
Show that a GLU MLP with $d_\text{ff}=\frac83d$ has the same parameter count and the same matmul FLOPs per
token as a plain MLP with $4d$. How many bytes of activations per token does each store for the
backward pass, in bf16, if only GEMM inputs and activation inputs are saved?
\end{exercise}

\begin{exercise}[NoPE and the causal mask]
Consider a single causal attention layer with all keys equal (so scores are uniform). Show that the
output at position $i$ is the mean of $v_0,\ldots,v_i$. Construct value vectors that let the layer output
a quantity strictly decreasing in $i$, demonstrating how position can be inferred without any encoding.
\end{exercise}

\begin{exercise}[z-loss gradient]
Compute $\partial(\log Z)^2/\partial z_j$. Decompose this gradient into a component along the all-ones
vector (which shifts every logit equally, changing $\log Z$ but not the predicted distribution) and an
orthogonal component. Show that the cross-entropy gradient has no all-ones component, and explain why, for
small $\alpha$, the z-loss controls $\log Z$ while barely perturbing the predictions.
\end{exercise}

\begin{exercise}[Norm fusion]
Using the napkin model of \cref{sec:archdesign:norm}, estimate for Qwen3-8B on an H100 the fraction of a
transformer block's forward time spent in unfused RMSNorms at 8192 tokens. Repeat for a fused
implementation that reads the input once and writes once inside the next GEMM's prologue.
\end{exercise}

\begin{exercise}[Aspect ratio and latency]
Two 1.3B-parameter models have $d/L=32$ and $d/L=128$. Ignoring embeddings, find $d$ and $L$ for each
(with $d_\text{ff}=4d$). Under a memory-bound decode where each layer adds a fixed 15\,$\mu$s of kernel launch
and synchronisation overhead plus bandwidth time, which model decodes faster on one GPU, and which one
splits more evenly over 4-way tensor parallelism?
\end{exercise}

\begin{exercise}[Design your own ablation]
Pick one modification not in \code{STANDARD\_ABLATIONS}, such as sandwich norm, parallel blocks, ALiBi,
QK-norm without z-loss or tied embeddings. Write down a hypothesis and the evidence that would refute it,
run it with two seeds and a 3-point LR sweep, and report a result with error bars.
\end{exercise}

% =====================================================================
\section{Lab 7: An architecture ablation workbench}\label{sec:archdesign:lab}
\begin{labbox}{7}{An architecture ablation workbench}
\labdir{lab07_ablations}
\labsection{Goals}
Implement the small components whose details decide whether a transformer trains well, verify each one
against a reference, and then use the configurable model in \code{llmsys/ablate.py} to run and analyse a
small ablation study.

\labsection{Setup}
\begin{lstlisting}[style=shell]
cd code
LAB_IMPL=starter  python -m pytest -q labs/lab07_ablations   # your work
LAB_IMPL=solution python -m pytest -q labs/lab07_ablations   # reference
pip install einops        # optional: to try the einops versions in the chapter
\end{lstlisting}

\labsection{Tasks}
\begin{enumerate}
\item \textbf{Einsum plumbing.} \code{einsum\_linear}, \code{split\_heads}, \code{merge\_heads}: write
  them with \code{torch.einsum}, then again with einops, and check both against \code{F.linear} and
  against a view-and-transpose implementation.
\item \textbf{Initialisation.} \code{fan\_avg\_std} and \code{trunc\_normal}, implemented by
  \emph{resampling} out-of-range draws rather than calling \code{nn.init.trunc\_normal\_}. Verify that
  the realised standard deviation is about 98.7\% of the target.
\item \textbf{MLP width.} \code{ffn\_hidden\_dim} for gated and plain MLPs with rounding to a multiple.
\item \textbf{Positions.} \code{sinusoidal\_positions} (check the relative-dot-product property in the
  test) and \code{alibi\_bias}.
\item \textbf{Norm placement.} \code{block\_forward} for pre, post, sandwich and parallel blocks, given
  arbitrary callables.
\item \textbf{Stability.} \code{z\_loss\_cross\_entropy} (numerically stable for logits of $\pm10^4$) and
  \code{softcap}.
\item \textbf{Accounting.} \code{parameter\_count(cfg)}: derive by hand the exact parameter count of
  \code{ArchLM(cfg)} for every option combination. The test compares it against the module.
\item \textbf{Reporting.} \code{summarize\_ablation}: average final losses across seeds, count
  divergences, and compute deltas against the baseline.
\end{enumerate}

\labsection{Measurements}
\begin{itemize}
\item Run \code{run.py --steps 400 --seeds 0 1} on CPU (a few minutes per variant) and plot the curves. Is
  any difference larger than your seed spread?
\item For pre-norm, post-norm and no-norm, sweep the learning rate over $\{10^{-3},3\cdot10^{-3},10^{-2},
  3\cdot10^{-2}\}$ and plot final loss against learning rate. Expect the normalised models to tolerate
  higher rates.
\item On a GPU, rerun the sweep on a real corpus with $d=512$, $L=8$ for about 5k steps. Which conclusions
  from the CPU sweep survive?
\end{itemize}

\labsection{Deliverables}
Passing tests; a learning-curve figure and an LR-sweep figure; a short table from
\code{summarize\_ablation}; and three sentences per variant explaining the result with the mechanisms in
this chapter.

\labsection{Stretch goals}
Add a \code{"peri"} placement that normalises only branch outputs (OLMo 2 style); add per-layer logging
of the maximum attention logit and show that QK-norm bounds it; reproduce \cref{fig:archdesign:depth}
for the actual \code{ArchLM} at initialisation for each norm placement.
\end{labbox}
